% Einheit 4 von 4 – Mit Größen rechnen im Sachzusammenhang (bank/einheiten/e4.jsonl, zusammenbau v0.1)
%% TODO Zweigzeile Teil 1: was hier gelernt wird (Satz in Schülersprache aus dem Einheitstitel) – Einheit 4
\einheitenkopf[e4]{Einheit 4 von 4 · Mit Größen rechnen im Sachzusammenhang}
\zweigzeile{ab Kl. 8 · P10}
\setcounter{aufgabe}{33}

%% TODO \verfahren{…}: Verfahrensüberschrift in Sachsprache fehlt in der Bank (2.3 b)
%% TODO Titel: Ich-kann-Satz fehlt in der Bank (Platzhalter: Kettenname) – Nr. 34
%% TODO Anweisung: ein Satz über der ersten Teilaufgabe fehlt in der Bank – Nr. 34
\begin{aufgabe}{Mit Größen rechnen im Sachzusammenhang}
\begin{teile}
\teil Von einer $1{,}2$ m langen Leiste werden $45$ cm abgesägt; gefragt ist der Rest in cm. Welche Angabe musst du vorher umrechnen? Kreuze an.\\ \kreuz{$1{,}2$ m}\\ \kreuz{$45$ cm}
\end{teile}
%% TODO Darstellung für schwach fehlt in der Bank (einheiten-e4-k1-s1-v1; Abschnitt „Für schwache Schüler“)
\swz{$1{,}4$ km $+ 650$ m – wie viel km?}{}
%% TODO Darstellung für schwach fehlt in der Bank (einheiten-e4-k1-s1-v2; Abschnitt „Für schwache Schüler“)
\swz{$2$ kg $+ 750$ g – wie viel kg?}{}
%% TODO Darstellung für schwach fehlt in der Bank (einheiten-e4-k1-s1-v3; Abschnitt „Für schwache Schüler“)
\swz{$3{,}2$ m $+ 45$ cm – wie viel m?}{}
%% TODO Darstellung für schwach fehlt in der Bank (einheiten-e4-k1-s1-v4; Abschnitt „Für schwache Schüler“)
\swz{$4{,}20$ € $+ 85$ ct – wie viel €?}{}
%% TODO Darstellung für schwach fehlt in der Bank (einheiten-e4-k1-s1-v5; Abschnitt „Für schwache Schüler“)
\swz{$1{,}5$ l $+ 330$ ml – wie viel l?}{}
\swfrage{Was bleibt gleich, was ändert sich?}
%% TODO Darstellung für schwach fehlt in der Bank (einheiten-e4-k1-s2-v1; Abschnitt „Für schwache Schüler“)
\swz{Ein Modellhaus ist $0{,}08$ m hoch. Gib die Höhe in einer handlicheren Einheit an.}{}
%% TODO Darstellung für schwach fehlt in der Bank (einheiten-e4-k1-s3-v1; Abschnitt „Für schwache Schüler“)
\swz{Ein Aluminiumwürfel hat ein Volumen von $30$ cm³; Aluminium hat die Dichte $2{,}7$ g/cm³. Masse?}{}
%% TODO Darstellung für schwach fehlt in der Bank (einheiten-e4-k1-s4-v1; Abschnitt „Für schwache Schüler“)
\swz{Ein Kupferteil wiegt $445$ g; Kupfer hat die Dichte $8{,}9$ g/cm³. Volumen?}{}
%% TODO Darstellung für schwach fehlt in der Bank (einheiten-e4-k1-s5-v1; Abschnitt „Für schwache Schüler“)
\swz{Ein Wassertank fasst $2$ m³ und wiegt leer $85$ kg. Er ist zu $30\,\%$ gefüllt; ein Kubikmeter Wasser wiegt eine Tonne. Gesamtmasse?}{}
%% TODO Darstellung für schwach fehlt in der Bank (einheiten-e4-k1-s6-v1; Abschnitt „Für schwache Schüler“)
\swz{Eine Wand hat $18$ m²; $1$ l Farbe reicht für $6$ m². Wie viele Liter braucht man?}{}
%% TODO Darstellung für schwach fehlt in der Bank (einheiten-e4-k1-s7-v1; Abschnitt „Für schwache Schüler“)
\swz{Zwei Wände, je $4$ m lang und $2{,}6$ m hoch, werden zweimal gestrichen; $1$ l Farbe reicht für $8$ m². Wie viele Liter?}{}
%% TODO Darstellung für schwach fehlt in der Bank (einheiten-e4-k1-s8-v1; Abschnitt „Für schwache Schüler“)
\swz{Man braucht $1{,}3$ l Farbe. Es gibt Dosen mit $0{,}75$ l, $1$ l und $2$ l. Welche einzelne Dose reicht mindestens?}{}
%% TODO Darstellung für schwach fehlt in der Bank (einheiten-e4-k1-s9-v1; Abschnitt „Für schwache Schüler“)
\swz[3]{Jonas rechnet seinen Schulweg: $1{,}6$ km $+ 750$ m $= 751{,}6$ m. Nenne den Fehler und rechne richtig.}{}
%% TODO Darstellung für schwach fehlt in der Bank (einheiten-e4-k1-s10-v1; Abschnitt „Für schwache Schüler“)
\swz[3]{Eine Messingkugel wiegt $3\,500$ g; Messing hat die Dichte $8{,}4$ g/cm³. Berechne das Volumen der Kugel. Runde auf eine Stelle nach dem Komma. \hfill (P10 2016 OS)}{}
%% TODO Darstellung für schwach fehlt in der Bank (einheiten-e4-k1-s11-v1; Abschnitt „Für schwache Schüler“)
\swz[3]{Ein Vogelhaus hat zwei Dachflächen zu je $0{,}18$ m² und vier Wände zu je $0{,}12$ m². Alles wird zweimal gestrichen; $1$ l Farbe reicht für $8$ m². Es gibt Dosen mit $125$ ml, $250$ ml und $750$ ml. Lea meint, die kleinste Dose reicht, Max will die mittlere kaufen. Wer hat recht? Begründe. \hfill (P10 2015 OS)}{}
\swz[3]{Ein Busunternehmen hat auf $120\,000$ km Mehrkosten von $744$ € für Reifen. In der Zeitung steht: „Jeder Kilometer wird dadurch etwa $6$ ct teurer.“ Die Rechnung dazu: $744$ € $= 744\,000$ ct; $744\,000 : 120\,000 = 6{,}2$ ct. Erkläre den Fehler und korrigiere die Behauptung. \hfill (P10 2014 OS)}{}
\end{aufgabe}

%% TODO \verfahren{…}: Verfahrensüberschrift in Sachsprache fehlt in der Bank (2.3 b)
%% TODO Titel: Ich-kann-Satz fehlt in der Bank (Platzhalter: Kettenname) – Nr. 35
%% TODO Anweisung: ein Satz über der ersten Teilaufgabe fehlt in der Bank – Nr. 35
\begin{aufgabe}{Kosten und Einnahmen, Sek II}
\begin{teile}
\teil Die Fläche ist in cm² gegeben, der Druckpreis in Euro pro m². Darf man direkt malnehmen? Kreuze an.\\ \kreuz{ja}\\ \kreuz{nein}
\end{teile}
%% TODO Darstellung für schwach fehlt in der Bank (einheiten-e4-k2-s1-v1; Abschnitt „Für schwache Schüler“)
\swz{Ein Kiosk verkauft $23$ Eis der Sorte A und $17$ der Sorte B, beide zu $2$ €. Einnahme?}{}
%% TODO Darstellung für schwach fehlt in der Bank (einheiten-e4-k2-s1-v2; Abschnitt „Für schwache Schüler“)
\swz{Ein Kino verkauft am Montag $48$ und am Dienstag $36$ Karten zu je $7$ €. Einnahme?}{}
%% TODO Darstellung für schwach fehlt in der Bank (einheiten-e4-k2-s1-v3; Abschnitt „Für schwache Schüler“)
\swz{Ein Bäcker verkauft Montag $120$, Dienstag $95$ und Mittwoch $105$ Brötchen zu je $0{,}45$ €. Einnahme?}{}
%% TODO Darstellung für schwach fehlt in der Bank (einheiten-e4-k2-s1-v4; Abschnitt „Für schwache Schüler“)
\swz{Ein Verein verkauft T-Shirts: $18$ in Größe S, $25$ in M und $21$ in L, alle zu $12$ €. Einnahme?}{}
%% TODO Darstellung für schwach fehlt in der Bank (einheiten-e4-k2-s1-v5; Abschnitt „Für schwache Schüler“)
\swz{Für ein Konzert gehen $136$ Karten im Vorverkauf und $89$ an der Abendkasse weg, alle zu $15{,}50$ €. Einnahme?}{}
\swfrage{Was bleibt gleich, was ändert sich?}
%% TODO Darstellung für schwach fehlt in der Bank (einheiten-e4-k2-s2-v1; Abschnitt „Für schwache Schüler“)
\swz{Ein Stand verkauft Stofftaschen: weiße $22$ und schwarze $17$ zu je $6$ €, bunte $15$, $9$ und $12$ zu je $4$ €. Einnahme?}{}
%% TODO Darstellung für schwach fehlt in der Bank (einheiten-e4-k2-s3-v1; Abschnitt „Für schwache Schüler“)
\swz{Ein Plakat hat $7\,500$ cm²; der Druck kostet $30$ € pro m². Kosten?}{}
\swz[3]{Aus einer $2{,}4$ m langen Latte werden zwei Leisten gesägt, $5$ und $3$ Längeneinheiten lang; eine Längeneinheit entspricht $25$ cm. Wie viel Prozent der Latte sind Verschnitt? \hfill (FHR 2022)}{}
\end{aufgabe}

%% TODO Titel: Ich-kann-Satz fehlt in der Bank (Platzhalter: Kettenname) – Nr. 36
%% TODO Anweisung: ein Satz über der ersten Teilaufgabe fehlt in der Bank – Nr. 36
\begin{aufgabe}{Dichte aus Masse und Volumen (Vorrat)}
%% TODO Darstellung für schwach fehlt in der Bank (einheiten-e4-k3-s1-v1; Abschnitt „Für schwache Schüler“)
\swz{Ein Stein wiegt $540$ g und hat ein Volumen von $200$ cm³. Dichte?}{}
\end{aufgabe}

%% TODO Titel: Ich-kann-Satz fehlt in der Bank (Platzhalter: Kettenname) – Nr. 37
%% TODO Anweisung: ein Satz über der ersten Teilaufgabe fehlt in der Bank – Nr. 37
\begin{aufgabe}{Euro und Cent in einer Rechnung}
%% TODO Darstellung für schwach fehlt in der Bank (einheiten-e4-k4-s1-v1; Abschnitt „Für schwache Schüler“)
\swz{Drei Brezeln zu je $85$ ct und ein Kakao zu $2{,}40$ € – zusammen?}{}
\end{aufgabe}

%% TODO Titel: Ich-kann-Satz fehlt in der Bank (Platzhalter: Kettenname) – Nr. 38
%% TODO Anweisung: ein Satz über der ersten Teilaufgabe fehlt in der Bank – Nr. 38
\begin{aufgabe}{Mit Größen rechnen im Sachzusammenhang – Fehler finden, Begründen, Anwendung}
\swz[3]{Tim schreibt: $7{,}25$ € $= 7\,250$ ct. Wo ist der Fehler?}{}
\swz[3]{Beim Sportfest werden $45$ Getränke zu $3$ € und $25$ zu $1$ € verkauft. Lina rechnet: Durchschnittspreis $2$ €, $70 \cdot 2 = 140$ €. Wo ist der Fehler?}{}
\swfrage{Ein Ergebnis lautet $0{,}007$ km. Begründe, warum man es besser in Metern angibt.}
\swfrage{Begründe, warum man bei $40$ Karten zu $8$ € und $5$ Karten zu $4$ € nicht mit dem Durchschnittspreis $6$ € rechnen darf.}
\swz[3]{Ein Fliesenleger belegt einen Badboden mit $6{,}3$ m². Eine Fliese ist $30$ cm mal $30$ cm groß. Wie viele Fliesen braucht er mindestens (ohne Verschnitt)?}{}
\end{aufgabe}

\uebersichtskasten{Erst umrechnen, dann rechnen: Alle Angaben in dieselbe Einheit bringen, rechnen, das Ergebnis in einer handlichen Einheit angeben. \\ Dichte: Masse gleich Volumen mal Dichte; umgestellt Volumen gleich Masse geteilt durch Dichte. ϱ = 7,8 g/cm³, m = 4000 g: V = 4000 : 7,8 ≈ 512,8 cm³. \\ Bedarf: Menge gleich Fläche geteilt durch Ergiebigkeit; jeder weitere Anstrich zählt noch einmal. 2,56 m² bei 10 m² je Liter: 0,256 l = 256 ml. \\ Fremde Rechnung prüfen: Zuerst die Umrechnung nachrechnen, dann die Rechnung – der Fehler steckt meist im Faktor. \\ Kosten und Einnahmen (Sek II): Fläche in Quadratmeter mal Preis je Quadratmeter; bei mehreren Preisstufen die Stückzahlen je Stufe addieren, mit dem Preis multiplizieren und die Teilbeträge summieren – nicht den Durchschnittspreis mit der Gesamtzahl. 54 Ketten zu 5 € und 46 Ketten zu 3 €: 54 · 5 + 46 · 3 = 408 €.}
